\documentclass[exercice,a4]{thedocument}%
\usepackage[french]{babel}%
\usepackage{amsmath}%
\startdatakeys
    \source{22GENMATNC1}
    \BO{2026}
    \cycle{4}
    \competencesDuSocle{CA,RA}
    \objectifsApprentissage{C4B3ca,C4B2eb,C4B3cb}
    \motsCles{construction,réciproque de Pythagore,trigonométrie,angle,Thalès,Pythagore}
\enddatakeys
\begin{document}%

\medskip

Un triangle MWB est tel que MB $= 7,5$cm; WB $= 4,5$~cm et MW $= 6$~cm.

\medskip

\begin{enumerate}
\item Sur la copie, construire le triangle MWB.
\item Montrer que le triangle MWB est rectangle en W.

\textbf{Rédiger la réponse en faisant apparaître les différentes étapes.}
\item Calculer la mesure de l'angle $\widehat{\text{BMW}}$. Arrondir le résultat au degré près.
\item
	\begin{enumerate}
		\item Placer le point F sur le segment [WB] tel que WF $= 3$~cm.
		\item Tracer la parallèle à (MB) passant par F. Elle coupe (MW) en E. Placer le point E.
		\item Calculer WE.

\textbf{Rédiger la réponse en faisant apparaître les différentes étapes.}
	\end{enumerate}
\item
	\begin{enumerate}
		\item Placer le point T sur la demi-droite [MW) de la figure précédente tel que MT $= 10$~cm.
		\item Tracer le segment [TB].
	\end{enumerate}
\item Calculer la longueur TE.

\textbf{Faire apparaître les différentes étapes du calcul.}
\end{enumerate}

\bigskip

\showthesource
\end{document}
